\subsection{\texorpdfstring{হেক্সাডেসিমাল থেকে $n$ ভিত্তি}{হেক্সাডেসিমাল থেকে n ভিত্তি}}
প্রথমে স্থানমান ব্যবহার করে দশমিকে আনো; তারপর লক্ষ্য ভিত্তিতে ভাগ ও গুণ করো। দুই ধাপের পূর্ণাংশ ও ভগ্নাংশ আলাদাভাবে মিলিয়ে নাও।
\par\noindent\begin{minipage}{\linewidth}
\convExample{১}{পূর্ণসংখ্যা}
$ (D3)_{16} $-কে ৪ ভিত্তিতে রূপান্তর করো।\par
\textbf{প্রথম ধাপ: দশমিক মান}\par
\par\noindent\begin{minipage}{\linewidth}\centering
\phantomsection\label{VC058}
\begin{chThreeDiagram}\begin{tikzpicture}[x=1mm,y=1mm]
\node[vc digit] (d0) at (0,0) {$ \mathrm{D} $};
\node[vcnum] (w0) at (0,-10) {$ 16^{1} $};
\node[vcnum] (v0) at (0,-25) {$ 208 $};
\draw[vc arrow] (d0.south)--(w0.north);
\draw[vc arrow] (w0.south)--(v0.north);
\node[vc digit] (d1) at (19,0) {$ 3 $};
\node[vcnum] (w1) at (19,-10) {$ 16^{0} $};
\node[vcnum] (v1) at (19,-25) {$ 3 $};
\draw[vc arrow] (d1.south)--(w1.north);
\draw[vc arrow] (w1.south)--(v1.north);
\end{tikzpicture}\end{chThreeDiagram}
\end{minipage}\par

\[\begin{aligned}V&=\mathrm{D}\times 16^{1}+3\times 16^{0}\\&=208+3=211.\end{aligned}\]
\textbf{দ্বিতীয় ধাপ: লক্ষ্য ভিত্তি}\par
\par\noindent\begin{minipage}{\linewidth}\centering
\phantomsection\label{VC059}
\begin{chThreeDiagram}\begin{tikzpicture}[x=1mm,y=1mm]
\path[use as bounding box] (-3,9) rectangle (72,-42);
\node[vcnote] at (15,7) {ভাজ্য/ভাগফল};
\node[vcnote] at (44,7) {ভাগশেষ};
\node[vcnum] at (0.0,0) {$ 4 $};
\node[vcnum] at (10.0,0) {$ 211 $};
\draw[vc line] (5.0,3)--(5.0,-3.5)--(50,-3.5);
\node[vcnum] (r0) at (44,-7) {$ 3 $};
\node[vcnum] at (1.5,-7) {$ 4 $};
\node[vcnum] at (11.5,-7) {$ 52 $};
\draw[vc line] (6.5,-4)--(6.5,-10.5)--(50,-10.5);
\node[vcnum] (r1) at (44,-14) {$ 0 $};
\node[vcnum] at (3.0,-14) {$ 4 $};
\node[vcnum] at (13.0,-14) {$ 13 $};
\draw[vc line] (8.0,-11)--(8.0,-17.5)--(50,-17.5);
\node[vcnum] (r2) at (44,-21) {$ 1 $};
\node[vcnum] at (4.5,-21) {$ 4 $};
\node[vcnum] at (14.5,-21) {$ 3 $};
\draw[vc line] (9.5,-18)--(9.5,-24.5)--(50,-24.5);
\node[vcnum] (r3) at (44,-28) {$ 3 $};
\node[vcnum] at (16.0,-28) {$0$};
\node[vcnote] at (61,-3) {LSB};
\node[vcnote] at (61,-33) {MSB};
\draw[vc arrow] (61,-28)--(61,-7);
\node[vcnote] at (31,-40) {নিচ থেকে ওপরে পড়ো};
\end{tikzpicture}\end{chThreeDiagram}
\end{minipage}\par
\[\therefore\ (D3)_{16}=(3103)_{4}.\]
\end{minipage}\par
\par\noindent\begin{minipage}{\linewidth}
\convExample{২}{শুধু ভগ্নাংশ}
$ (0.A4)_{16} $-কে ৪ ভিত্তিতে রূপান্তর করো।\par
\textbf{প্রথম ধাপ: দশমিক মান}\par
\par\noindent\begin{minipage}{\linewidth}\centering
\phantomsection\label{VC060}
\begin{chThreeDiagram}\begin{tikzpicture}[x=1mm,y=1mm]
\node[vc digit] (d0) at (0,0) {$ 0 $};
\node[vcnum] (w0) at (0,-10) {$ 16^{0} $};
\node[vcnum] (v0) at (0,-25) {$ 0 $};
\draw[vc arrow] (d0.south)--(w0.north);
\draw[vc arrow] (w0.south)--(v0.north);
\node[vc digit] (d1) at (27,0) {$ \mathrm{A} $};
\node[vcnum] (w1) at (27,-10) {$ 16^{-1} $};
\node[vcnum] (v1) at (27,-25) {$ \frac{5}{8} $};
\draw[vc arrow] (d1.south)--(w1.north);
\draw[vc arrow] (w1.south)--(v1.north);
\node[vc digit] (d2) at (46,0) {$ 4 $};
\node[vcnum] (w2) at (46,-10) {$ 16^{-2} $};
\node[vcnum] (v2) at (46,-25) {$ \frac{1}{64} $};
\draw[vc arrow] (d2.south)--(w2.north);
\draw[vc arrow] (w2.south)--(v2.north);
\node[vc digit] at (13.5,0) {$.$};
\end{tikzpicture}\end{chThreeDiagram}
\end{minipage}\par

\[\begin{aligned}V&=0\times 16^{0}+\mathrm{A}\times 16^{-1}+4\times 16^{-2}\\&=\frac{5}{8}+\frac{1}{64}=0.640625.\end{aligned}\]
\textbf{দ্বিতীয় ধাপ: লক্ষ্য ভিত্তি}\par
\par\noindent\begin{minipage}{\linewidth}\centering
\phantomsection\label{VC061}
\begin{chThreeDiagram}\begin{tikzpicture}[x=1mm,y=1mm]
\path[use as bounding box] (-9,9) rectangle (76,-51);
\node[vcnote] at (14,7) {প্রাপ্ত অঙ্ক};
\node[vcnote] at (45,7) {ভগ্নাংশ};
\node[vcnum] at (45,0) {$ 0.640625 $};
\draw[vc line] (26,3)--(26,-39);
\node[vcnum] (m0) at (45,-5) {$\times 4$};
\draw[vc line] (7,-8)--(65,-8);
\node[vcnum] at (14,-12) {$ 2 $};
\node[vcnum] (f0) at (45,-12) {$ 0.5625 $};
\node[vcnum] (m1) at (45,-17) {$\times 4$};
\draw[vc line] (7,-20)--(65,-20);
\node[vcnum] at (14,-24) {$ 2 $};
\node[vcnum] (f1) at (45,-24) {$ 0.25 $};
\node[vcnum] (m2) at (45,-29) {$\times 4$};
\draw[vc line] (7,-32)--(65,-32);
\node[vcnum] at (14,-36) {$ 1 $};
\node[vcnum] (f2) at (45,-36) {$ 0 $};
\draw[vc arrow] (f0.east)--(72,-12)--(72,-17)--(m1.east);
\draw[vc arrow] (f1.east)--(72,-24)--(72,-29)--(m2.east);
\node[vcnote] at (0,-8) {MSB};
\node[vcnote] at (0,-41) {LSB};
\draw[vc arrow] (0,-12)--(0,-36);
\node[vcnote] at (34,-48) {ওপর থেকে নিচে পড়ো};
\end{tikzpicture}\end{chThreeDiagram}
\end{minipage}\par
\[\therefore\ (0.A4)_{16}=(0.221)_{4}.\]
\end{minipage}\par
\par\noindent\begin{minipage}{\linewidth}
\convExample{৩}{মিশ্র সংখ্যা}
$ (2F.8)_{16} $-কে ৪ ভিত্তিতে রূপান্তর করো।\par
\textbf{প্রথম ধাপ: দশমিক মান}\par
\par\noindent\begin{minipage}{\linewidth}\centering
\phantomsection\label{VC062}
\begin{chThreeDiagram}\begin{tikzpicture}[x=1mm,y=1mm]
\node[vc digit] (d0) at (0,0) {$ 2 $};
\node[vcnum] (w0) at (0,-10) {$ 16^{1} $};
\node[vcnum] (v0) at (0,-25) {$ 32 $};
\draw[vc arrow] (d0.south)--(w0.north);
\draw[vc arrow] (w0.south)--(v0.north);
\node[vc digit] (d1) at (19,0) {$ \mathrm{F} $};
\node[vcnum] (w1) at (19,-10) {$ 16^{0} $};
\node[vcnum] (v1) at (19,-25) {$ 15 $};
\draw[vc arrow] (d1.south)--(w1.north);
\draw[vc arrow] (w1.south)--(v1.north);
\node[vc digit] (d2) at (46,0) {$ 8 $};
\node[vcnum] (w2) at (46,-10) {$ 16^{-1} $};
\node[vcnum] (v2) at (46,-25) {$ \frac{1}{2} $};
\draw[vc arrow] (d2.south)--(w2.north);
\draw[vc arrow] (w2.south)--(v2.north);
\node[vc digit] at (32.5,0) {$.$};
\end{tikzpicture}\end{chThreeDiagram}
\end{minipage}\par

\[\begin{aligned}V&=2\times 16^{1}+\mathrm{F}\times 16^{0}+8\times 16^{-1}\\&=32+15+\frac{1}{2}=47.5.\end{aligned}\]
\textbf{দ্বিতীয় ধাপ: লক্ষ্য ভিত্তি}\par
\begin{minipage}[t]{.48\linewidth}\vspace{0pt}\centering\textbf{পূর্ণাংশ}\par\vspace{1mm}
\par\noindent\begin{minipage}{\linewidth}\centering
\phantomsection\label{VC063}
\begin{chThreeDiagram}\begin{tikzpicture}[x=1mm,y=1mm]
\path[use as bounding box] (-3,9) rectangle (72,-35);
\node[vcnote] at (15,7) {ভাজ্য/ভাগফল};
\node[vcnote] at (44,7) {ভাগশেষ};
\node[vcnum] at (0.0,0) {$ 4 $};
\node[vcnum] at (10.0,0) {$ 47 $};
\draw[vc line] (5.0,3)--(5.0,-3.5)--(50,-3.5);
\node[vcnum] (r0) at (44,-7) {$ 3 $};
\node[vcnum] at (1.5,-7) {$ 4 $};
\node[vcnum] at (11.5,-7) {$ 11 $};
\draw[vc line] (6.5,-4)--(6.5,-10.5)--(50,-10.5);
\node[vcnum] (r1) at (44,-14) {$ 3 $};
\node[vcnum] at (3.0,-14) {$ 4 $};
\node[vcnum] at (13.0,-14) {$ 2 $};
\draw[vc line] (8.0,-11)--(8.0,-17.5)--(50,-17.5);
\node[vcnum] (r2) at (44,-21) {$ 2 $};
\node[vcnum] at (14.5,-21) {$0$};
\node[vcnote] at (61,-3) {LSB};
\node[vcnote] at (61,-26) {MSB};
\draw[vc arrow] (61,-21)--(61,-7);
\node[vcnote] at (31,-33) {নিচ থেকে ওপরে পড়ো};
\end{tikzpicture}\end{chThreeDiagram}
\end{minipage}\par

$ (233)_{4} $
\end{minipage}\hfill
\begin{minipage}[t]{.48\linewidth}\vspace{0pt}\centering\textbf{ভগ্নাংশ}\par\vspace{1mm}
\par\noindent\begin{minipage}{\linewidth}\centering
\phantomsection\label{VC064}
\begin{chThreeDiagram}\begin{tikzpicture}[x=1mm,y=1mm]
\path[use as bounding box] (-9,9) rectangle (76,-27);
\node[vcnote] at (14,7) {প্রাপ্ত অঙ্ক};
\node[vcnote] at (45,7) {ভগ্নাংশ};
\node[vcnum] at (45,0) {$ 0.5 $};
\draw[vc line] (26,3)--(26,-15);
\node[vcnum] (m0) at (45,-5) {$\times 4$};
\draw[vc line] (7,-8)--(65,-8);
\node[vcnum] at (14,-12) {$ 2 $};
\node[vcnum] (f0) at (45,-12) {$ 0 $};
\node[vcnote] at (0,-10) {MSB=LSB};
\draw[vc arrow] (0,-14)--(0,-20);
\node[vcnote] at (34,-24) {ওপর থেকে নিচে পড়ো};
\end{tikzpicture}\end{chThreeDiagram}
\end{minipage}\par

$ (0.2)_{4} $
\end{minipage}\par
\[\therefore\ (2F.8)_{16}=(233.2)_{4}.\]
\end{minipage}\par
